\section{Related Work}
\label{sec:related_work}

\emph{Drafting scope.} This section will connect the observed results to five
specific bodies of work: classical MSI peak picking, msiPL, S3PL, spatial scale
and tissue-boundary mixing, and the limits of attribution as biological
evidence. It will be completed through a targeted literature pass rather than
by reusing the proposal's broad survey unchanged.

msiPL established that a VAE can learn an unsupervised spectral representation
and support model-derived peak selection~\citep{abdelmoula2021peak}. S3PL
subsequently incorporated spatial-spectral patches into a convolutional
architecture~\citep{weigand2026spatial}. The present study does not treat these
as architecture-matched baselines: their objectives, normalisation, training
protocols and selectors differ. Instead, it reproduces their executable
pipelines where possible and uses matched final peak counts to prevent a method
from improving a spatial score merely by returning more bins.

Integrated Gradients attributes the change in a differentiable target relative
to a specified baseline while satisfying a completeness property under its
standard assumptions~\citep{sundararajan2017axiomatic}. Here, the target is a
GMM posterior fitted to frozen latent means. This explains the evaluated model's
latent separation; it does not identify molecular compounds or establish that a
ranked bin is a validated biomarker.

